%%%%%%%%%%%%%%%%%%%%%%%%%%%%%%%%%%%%%%%%%%%%%%%%%%%%%%%%%%%%%%%%%%%%%%
% Lớp 11 | Quan hệ vuông góc trong không gian | Hai mặt phẳng vuông góc
% Dạng 1H8?4-5: Nhận dạng và tính toán liên quan các hình thông dụng
% Số câu: 5  (sinh tự động bởi tach_cau.py; từ nay thầy sửa trực tiếp file này)
%%%%%%%%%%%%%%%%%%%%%%%%%%%%%%%%%%%%%%%%%%%%%%%%%%%%%%%%%%%%%%%%%%%%%%

%%%%%%%%%%%%%%%% TRẮC NGHIỆM 4 PHƯƠNG ÁN %%%%%%%%%%%%%%%%

% nguồn: Ok_TracNghiemLop11_Lan1
\begin{ex}%[1H8V4-5]%[ma:00c8e104]
	Cho hình chóp $S.ABCD$ có $SA\perp (ABCD)$ và đáy  $ABCD$ là hình vuông. Gọi $M$, $N$ lần lượt là trung điểm của $SB$, $SD$. Mặt phẳng $(P)$ chứa $MN$ và vuông góc với mặt phẳng $(ABCD)$. Mặt phẳng $(P)$ cắt $AC$ tại $K$. Tỷ số $\dfrac{KC}{KA}$ bằng
	\choice 
		{\True  $3$}
		{$2$} 
		{$4$} 
		{$\dfrac{3}{2}$} 
	\loigiai{ 
	\immini{Gọi $P$, $Q$ lần lượt là  trung điểm của $AD$, $AB$.\\
		Ta có $MN\parallel SA$, $SA\perp (ABCD)$\\$\Rightarrow MQ\perp (ABCD)\Rightarrow (MNPQ)\perp (ABCD)$.\\ 
		Vậy $(P)\equiv (MNPQ)$.\\ 
		 Gọi $K=PQ\cap AC\Rightarrow K=(P)\cap AC$.\\
		Theo tính chất đường trung bình ta  có\\
		$\dfrac{KA}{AC}=\dfrac{KA}{2AO}=\dfrac{1}{4}\Rightarrow \dfrac{KC}{KA}=3$.} 
		{\begin{tikzpicture}[scale=0.9,line join=round,font=\footnotesize,line cap=round]
		\coordinate (A) at (0,0); 
		\coordinate  (B) at (-1.5,-2);
		\coordinate (D) at  (4,0);
		\coordinate (C) at ($(B)+(D)-(A)$); 
		\coordinate (S) at ($(A)+(0,4)$); 
		\path  ($(A)!0.5!(C)$) coordinate (O);
		\path ($(A)!0.5!(B)$)  coordinate (Q);
		\path ($(A)!0.5!(D)$) coordinate (P); 
		\path ($(S)!0.5!(B)$)  coordinate (M);
		\path ($(S)!0.5!(D)$) coordinate  (N);
		\path ($(A)!0.5!(O)$)  coordinate (K);
		 \draw (S)--(B)
			(S)--(C) 
			 (S)--(D)
			 (B)--(C)--(D);
		 \draw[dashed](S)--(A)
			 (A)--(B)
			 (A)--(D)
			(M)--(N)--(P)--(Q)--cycle  (A)--(C)
			(B)--(D); 
		\pic[draw,angle radius=3mm] {right angle = S--A--D} pic[draw,angle  radius=3mm] {right angle = S--A--B};
		\foreach  \i/\g in {S/90,A/-90,B/-90,C/-90,D/0,M/180,N/0,P/45,Q/180,K/90,O/-90}{\draw[fill](\i) circle (1pt)
			 ($(\i)+(\g:3mm)$) node[scale=1]{$\i$};}
	\end{tikzpicture}
	 }}
\end{ex}

% nguồn: Ok_TracNghiemLop11_Lan1
\begin{ex}%[1H8V4-5]%[ma:3877a22b]
	 Cho hình chóp $S.ABCD$ có đáy $ABCD$ là hình vuông cạnh $a$. Tam giác $SAB$ là tam giác vuông tại $S$ và nằm trong mặt phẳng vuông góc với mặt phẳng chứa đáy. Biết $I$ là một điểm trong không gian cách đều các điểm $A, B,C, D$ và $S$. Tính độ dài đoạn thẳng $IS$.
	 \choice
		{$IS=a$} 
		{$IS=a\sqrt{2}$} 
		{\True  $IS=\dfrac{a\sqrt{2}}{2}$}
		{$IS=\dfrac{a}{2}$} 
	\loigiai{ 
	 \immini{Gọi $M$ là trung điểm $AB$, $I$ là tâm hình vuông vuông $ABCD$. Ta có $IM$ là trục đường tròn ngoại tiếp $\Delta SAB\Rightarrow IA=IB=IS$. Lại vì $I$ là tâm của $ABCD$ nên có $IA=IB=IC=ID$. Từ đó suy ra $I$ là điểm cách đều các điểm $A,B,C,D,S$. Ta có $IS=\sqrt{IM^2+SM^2}=\sqrt{\dfrac{AB^2}{4}+\dfrac{AB^2}{4}}=\sqrt{\dfrac{a^2}{4}+\dfrac{a^2}{4}}=\dfrac{a\sqrt{2}}{2}$.}
		{\begin{tikzpicture}[>=stealth,line join=round,line cap=round,font=\footnotesize,scale=2]
		 \draw[scale=2.5] (0,0) coordinate (A)
			(-120:0.5)  coordinate (B)--([turn]120:1) coordinate (C)--(1,0) coordinate (D)--($(A)+(100:0.6)$) coordinate (S)
			 (B)--(S)--(C);
		\draw[dashed,black,scale=2.5]  (B)--(A)--(D)
			 (S)--(A);
		\coordinate (I) at ($(B)!0.5!(D)$); 
		\coordinate  (M) at ($(A)!0.5!(B)$);
		\coordinate (P)  at ($(S)!2/3!(A)$);
		 \draw[dashed] (A)--(C)
			 (B)--(D)
			 (S)--(M)--(I)--(S);
		\foreach  \diem/\vitrin in {M/below right,I/below,S/above,A/above right,B/below left,C/below right,D/right} \fill (\diem) circle (1pt) node[\vitrin]{$\diem$};
	\end{tikzpicture}
		} 
	} 
\end{ex}


%%%%%%%%%%%%%%%% ĐÚNG/SAI %%%%%%%%%%%%%%%%

% nguồn: Ok_dungsai_lop11_Lan1
\begin{ex}%[1H8H4-5]%[ma:7e17a982]
 Cho hình chóp $S.ABCD$  có đáy $ABCD$ là hình thoi tâm $O$, đường thẳng $SO$ vuông góc với mặt phẳng $(ABCD)$. Biết $AB=SB=a$, $SO=\dfrac{a\sqrt{6}}{3}$. Khi đó
 \choiceTF 
  {\True $AC\perp  (SBD)$}
  {$\left((SAC),(SBD)\right)=60^\circ$} 
   {\True $BD=\dfrac{2a\sqrt{3}}{3}$}
  {\True $(SAB)\perp (SAD)$} 
 \loigiai{ 
  \begin{center}
   \begin{tikzpicture}[scale=1,line join=round,font=\footnotesize,line cap=round,thick]
     \coordinate (A) at  (0,0);
     \coordinate (B) at  (-2.5,-2);
      \coordinate (D) at (5,0);
     \coordinate  (C) at ($(B)+(D)-(A)$);
     \coordinate (O)  at ($(A)!0.5!(C)$);
     \coordinate (S) at  ($(O)+(0,5)$);
     \coordinate (M)  at ($(A)!0.5!(S)$);
     \draw  (S)--(B)
      (S)--(C) 
       (B)--(C)
      (C)--(D) 
      (S)--(D); 
      \draw[dashed,thin](A)--(C)
      (A)--(D) 
      (S)--(O) 
       (B)--(D)
       (S)--(A)
       (A)--(B)
      (B)--(M)--(D); 
     \pic[draw,thin,angle radius=2mm] {right angle =  S--O--D} pic[draw,thin,angle radius=2mm] {right angle = S--O--A};
     \foreach  \i/\g in {S/90,A/180,B/-90,C/-90,D/0,O/-90,M/45}{\draw[fill=black](\i) circle (1pt)
      ($(\i)+(\g:3mm)$) node[scale=1]{$\i$};} 
    \end{tikzpicture}
  \end{center}
  \begin{itemchoice}
    \itemch{ Đúng.}\\
   Ta có $\heva{&AC\perp BD\\&AC\perp SO}\Rightarrow  AC\perp (SBD)$.
    \itemch{ Sai.}\\
   Do $\heva{&AC\perp (SBD)\\&AC\subset  (SAC)}\Rightarrow (SAC)\perp (SBD)$.
   \itemch{ Đúng.}\\ 
   Xét $\triangle SOB$ vuông tại $O$,  ta có $OB^2=\sqrt{SB^2-SO^2}=\sqrt{a^2-\dfrac{2a^2}{3}}=\dfrac{a\sqrt{3}}{3}$\\$\Rightarrow BD=2OB=\dfrac{2a\sqrt{3}}{3}$.
    \itemch{ Đúng.}\\
   Gọi $M$ là trung điểm của $SA$.\\ 
   Ta có $\triangle  ABS$ cân tại $B$ và $\triangle ADS$ cân tại $D$ nên $BM$ và $DM$ cùng vuông góc với $SA$.\\
   Khi đó  góc giữa hai mặt phẳng $(SAB)$ và $(SAD)$ bằng hoặc bù với $\widehat{ BMD}$.\\
   Xét $\triangle ABO$ vuông tại $O$, ta có $AO=\sqrt{AB^2-OB^2}=\sqrt{a^2-\dfrac{a^2}{3}}=\dfrac{a\sqrt{6}}{3}=SO$.\\$\Rightarrow\triangle SAO$ vuông cân tại $O$ nên $SA=SO\cdot \sqrt{2}=\dfrac{2a\sqrt{3}}{3}$.\\ 
   Do đó $OM=\dfrac{1}{2}SA=\dfrac{a\sqrt{3}}{3}=\dfrac{1}{2}BD\Rightarrow \triangle BMD$ vuông  cân tại $M\Rightarrow \widehat{BMD}=90^\circ$.\\
   Vậy  $(SAB)\perp (SAD)$.
  \end{itemchoice}
  }
\end{ex}

% nguồn: Ok_dungsai_lop11_Lan1
\begin{ex}%[1H8H4-5]%[ma:da95f5f2]
 Cho lăng trụ đứng $ABCD.A'B'C'D'$ có đáy là  tam giác vuông tại $A$, biết $AB=a$, $AC=a\sqrt{3}$ và $\left((ACB'),(ABC)\right)=60^\circ $. Khi đó
 \choiceTF 
   {\True $A'A\perp (ABC)$}
   {$\left((ACB'),(ABB'A')\right)=60^\circ$}
  {\True $\left((ACC'A'),(BCCB')\right)=30^\circ$} 
  {\True Tổng diện tích ba  mặt bên của hình lăng trụ đã cho bằng $(3\sqrt{3}+3)a^2$}
 \loigiai{ 
  \begin{center}
   \begin{tikzpicture}[scale=1,line join=round,font=\footnotesize,line cap=round,thick]
      \def\a{4}
     \def\h{4.5} 
     \path (0:0) coordinate  (A)
       ++(0:\a) coordinate (C)
       ++(-150:3*\a/4) coordinate (B)
      ($(A)+(90:\h)$) coordinate  (A')
      ($(B)+(90:\h)$)  coordinate (B')
       ($(C)+(90:\h)$) coordinate (C');
      \draw[dashed,thick] (A)--(C);
     \draw[thick](A)--(B')--(C) 
      (C)--(C') 
       (B)--(B')
      (A)--(A') 
      (A)--(B)--(C) 
       (A')--(B')--(C')--cycle;
     \foreach  \x/\g in {A/180,B/-45,C/0,A'/180,B'/90,C'/0}
      \fill[black] (\x) circle (1pt)
       ($(\g:4mm)+(\x)$) node {$\x$};
      \pic[draw,thin,angle radius=3mm] {right angle = B--A--C};
     \pic[draw,thin,angle radius=3mm] {right angle  = B'--A'--C'};
     \path ($(A)!0.5!(B)$)  node [left]{$a$};
     \path ($(A)!0.5!(C)$) node [above]{$a\sqrt{3}$}; 
    \end{tikzpicture}
  \end{center}
  \begin{itemchoice}
   \itemch{ Đúng.}\\ 
   Vì $ABC.A'B'C'$ là lăng trụ  đứng nên $A'A\perp (ABC)$.
   \itemch{  Sai.}\\
    Ta có $A'A\perp (ABC)\Rightarrow A'A\perp AC$.\\
   Mặt khác  $AB\perp AC$,
    do đó $AC\perp (ABB'A')$.\\
   Mà $AC\subset (ACB')$  nên $(ACB')\perp (ABB'A')$.
   \itemch{  Đúng.}\\
   Ta có $\heva{&CC'=(ACC'A')\cap (BCC'B')\\&AC\perp CC',BC\perp  CC'\\&AC\subset (ACC'A'), BC\subset (BCC'B')}$\\
   $\Rightarrow  \left((ACC'A'),(BCC'B')\right)=(AC,BC)=\widehat{ACB}$.\\
   Tam giác $ABC$ vuông  tại $A$ có $\tan\widehat{ACB}=\dfrac{AB}{AC}=\dfrac{a}{a\sqrt{3}}=\dfrac{1}{\sqrt{3}}\Rightarrow \widehat{ACB}=30^\circ$.
    \itemch{ Đúng.}\\
   Ta có $\heva{&AC=(ACB')\cap (ABC)\\&AB\perp  AC, AB'\perp AC\\&AB\subset (ABC), AB'\subset (ACB')}$\\
   $\Rightarrow \left((ACB'),(ABC) \right)=(AB',AB)=\widehat{BAB'}=60^\circ$.\\ 
   Tam giác $ABB'$ vuông tại $B$  có $BB'=AB\cdot \tan 60^\circ=a\sqrt{3}$.\\
   Tam giác $ABC$ vuông tại $A$ có  $BC=\sqrt{AB^2+AC^2}=2a$.\\
    Tổng diện tích các mặt bên của hình trụ: $a\cdot a\sqrt{3}+2a\cdot a\sqrt{3}+a\sqrt{3}\cdot a\sqrt{3}=(3\sqrt{3}+3)a^2$.
  \end{itemchoice}
  }
\end{ex}


%%%%%%%%%%%%%%%% TRẢ LỜI NGẮN %%%%%%%%%%%%%%%%

% nguồn: Ok_TraLoiNgan_Lop11_Lan1
\begin{ex}%[1H8V4-5]%[ma:1cd59e3b]
	Cho hình chóp cụt tứ giác đều có cạnh đáy nhỏ là $a$, cạnh đáy lớn là $2a$ và chiều cao là $3a$. Độ dài cạnh bên của  hình chóp cụt bằng $m\cdot a$. Giá trị $m$ là bao nhiêu? (kết quả làm tròn đến hàng phần trăm)
	\shortans{3,08} 
	 \loigiai{
	 \immini{
		Kẻ  $C'H\perp OC$.\\
		Ta có  $$HC=OC-OH=\dfrac{2a\sqrt{2}}{2}-\dfrac{a\sqrt{2}}{2}=\dfrac{a\sqrt{2}}{2}.$$
		Trong tam giác $CC'H$ vuông, ta có $$CC'=\sqrt{C'H^2+HC^2}=\sqrt{9a^2+\dfrac{a^2}{2}}=\dfrac{\sqrt{38}}{2}a.$$\\Vậy $m=\dfrac{\sqrt{38}}{2}\approx 3{,}08$.}{\begin{tikzpicture}[scale=1,line join=round,font=\footnotesize,line cap=round]
		\def\a{5.5} 
		 \def\h{6}
		 \def\tiso{0.5}
		\path  (0:0) coordinate (A)
			++(0:\a) coordinate  (B)
			($(A)+(-130:\a/2)$) coordinate (D) 
			($(B)+(D)-(A)$)  coordinate (Ct)
			($(D)!1!(Ct)$) coordinate  (C)
			($(A)!0.5!(C)$) coordinate (H); 
		\coordinate (S) at  ($(H)+(90:\h)$);
		\path \foreach \x in  {A,B,C,D}{($(\x)!\tiso!(S)$) coordinate (\x')};
		\path ($(A)!0.5!(C)$) coordinate (O); 
		\path ($(A')!0.5!(C')$) coordinate  (O');
		\path  ($(O)!0.5!(C)$) coordinate (H);
		\draw[dashed] (B)--(A)--(D) 
			 (A)--(A')
			 (A)--(C)
			(B)--(D) 
			(O)--(O') 
			(C')--(H); 
		 \draw (B)--(C)--(D)
			(D')--(C')--(B')--(A')--cycle 
			(B)--(B') 
			 (C)--(C')
			(D)--(D') 
			(A')--(C') 
			 (B')--(D');
		\foreach \x/\g in  {A/135,B/45,C/-45,D/-135,A'/135,B'/45,C'/-35,D'/180,O/-90,O'/90,H/-90}
		\fill[black]  (\x) circle (1pt)
			($(\g:4mm)+(\x)$) node {$\x$}; 
	\end{tikzpicture}
		 }
	} 
\end{ex}
